\chapter{Notation and Conventions}

\begin{table}[htbp]
\centering
\caption{Names used throughout the book.}
\begin{tabular}{@{}lp{8.5cm}@{}}
\toprule
Name & Meaning \\
\midrule
XLEN & Width of an integer register, in bits. 32 or 64. \\
Datapath & Width of the integer ALU and of the data bus, in bits. \\
ELEN & Widest vector element the implementation supports. \\
VLEN & Bits in one vector register. A power of two, at least ELEN. \\
SEW & Selected element width for the current configuration. \\
EEW & Effective element width of one instruction. \\
VL & Number of elements the current instruction processes. \\
VLMAX & Maximum VL for the current SEW and LMUL. \\
LMUL & Register grouping factor. \\
SVM-8 & BX-8 profile: ELEN 8, VLEN 32, EEW 8. \\
SVM-16 & BX-16 profile: ELEN 16, VLEN 128, EEW 8 and 16. \\
\bottomrule
\end{tabular}
\end{table}

Instructions are written in the usual assembly form, \texttt{vadd.vv v1, v2, v3}. Configuration is \texttt{vsetvl} or \texttt{vsetvli}. The thirty-two vector registers are \texttt{v0} through \texttt{v31}. The mask register is \texttt{v0}.

A \emph{profile} is a named subset of RVV behaviour that this book specifies and that the specification does not ratify under that name. A \emph{ratified extension} is a string the RISC-V manuals define: Zve32f, Zve64d, or V.

Bluespec fragments are illustrative. They show interfaces and rules, not a vendor core. Module names use the prefix \texttt{Bx}.

The words MSX and Amiga, in this book, name a peripheral and bus contract. They do not name the original microprocessor. BX-8 is MSX-compatible in that contract and can replace an MSX in a rack or on a desk. BX-16 is Amiga-compatible in that contract and can replace an Amiga in the same sense.
